\section{Overlap presentations and the sharp threshold}\label{sec:stab:conjugacy}

The local fibers become a dynamical problem only after overlap is imposed.  In
this section we pass from the finite-window map $\Fold_m$ to a sliding-block
factor, prove the finite-state presentation of its image, and then prove the
sharp three-window reconstruction threshold.  The terminology for sofic shifts
and labeled graph presentations is standard; see
\cite{LindMarcus1995,Kitchens1998,Nasu1995}.

\begin{definition}[Sliding stabilized-window process]\label{def:Phi}
Fix $m\ge 2$.  The two-sided sliding stabilized-window map is
\[
\Phi_m:\{0,1\}^{\ZZ}\longrightarrow (X_m)^{\ZZ},
\qquad
(\Phi_m(a))_t:=\Fold_m(a_t,a_{t+1},\ldots,a_{t+m-1}).
\]
Its image is denoted
\[
Y_m:=\Phi_m(\{0,1\}^{\ZZ})\subset (X_m)^{\ZZ}.
\]
For $n\ge 1$ the associated finite-block map is
\[
\Phi_{m,n}:\{0,1\}^{n+m-1}\longrightarrow X_m^n,
\qquad
\Phi_{m,n}(a_0\ldots a_{n+m-2})
:=
\bigl(\Fold_m(a_t\ldots a_{t+m-1})\bigr)_{t=0}^{n-1}.
\]
We write residues of stabilized windows as
\[
r_t:=\Val_m\bigl((\Phi_m(a))_t\bigr).
\]
Equivalently,
$r_t\equiv\Val_m(a_t\ldots a_{t+m-1})\pmod {F_{m+2}}$.
\end{definition}

\begin{definition}[The overlap graph]\label{def:overlap-graph}
For $m\ge2$, let $\mathcal G_m$ be the labeled directed graph with vertex set
\[
V_m:=\{0,1\}^{m-1}.
\]
For every vertex $u=(u_1,\ldots,u_{m-1})\in V_m$ and every
$c\in\{0,1\}$ there is an edge
\[
u\longrightarrow (u_2,\ldots,u_{m-1},c)
\]
with label
\[
\ell(u,c):=\Fold_m(u_1,\ldots,u_{m-1},c)\in X_m.
\]
Thus a path in $\mathcal G_m$ records the successive length-$(m-1)$ overlaps
of a raw binary sequence, while the edge label records the stabilized
length-$m$ window.
\end{definition}

\paragraph{Finite-state lift and bounded multiplicity.}\label{thm:finite-state-lift}
For every $m\ge2$ the image shift $Y_m$ is presented by the labeled graph
$\mathcal G_m$.  Moreover $\mathcal G_m$ is both right-resolving and
left-resolving, and every point of $Y_m$ has at most $2^{m-1}$ preimages under
$\Phi_m$.

\begin{proof}
A bi-infinite raw sequence $a\in\{0,1\}^{\ZZ}$ determines the path whose
vertex at time $t$ is
\[
u_t=(a_t,\ldots,a_{t+m-2})\in V_m.
\]
The edge from $u_t$ to $u_{t+1}$ is obtained by appending $a_{t+m-1}$, and its
label is
\[
\ell(u_t,a_{t+m-1})=\Fold_m(a_t,\ldots,a_{t+m-1})=(\Phi_m(a))_t.
\]
Thus every point of $Y_m$ is read along a path in $\mathcal G_m$.  Conversely,
any bi-infinite path in $\mathcal G_m$ has overlapping vertices
\[
u_t=(a_t,\ldots,a_{t+m-2})
\]
for a unique raw sequence $a\in\{0,1\}^{\ZZ}$, and its label sequence is
$\Phi_m(a)$.  Hence $\mathcal G_m$ presents exactly $Y_m$.

To prove right-resolvingness, fix $u\in V_m$ and suppose the two outgoing
edges corresponding to $c,c'\in\{0,1\}$ have the same label.  Since
$\Val_m|_{X_m}$ is injective, their labels have the same residue modulo
$F_{m+2}$.  The two raw windows differ only in the last coordinate, so
\[
(c-c')F_{m+1}\equiv0\pmod {F_{m+2}}.
\]
Because $\gcd(F_{m+1},F_{m+2})=1$, this forces $c=c'$.  Thus outgoing edges
from the same vertex have distinct labels.

The left-resolving proof is the same in reverse.  Fix a terminal vertex
$v=(v_1,\ldots,v_{m-1})$ and suppose the two incoming edges obtained from
first bits $b,b'\in\{0,1\}$ have the same label.  The raw windows differ only
in their first coordinate, whose weight is $F_2=1$, hence
\[
b-b'\equiv0\pmod {F_{m+2}}.
\]
Since $F_{m+2}\ge3$, this gives $b=b'$.  Thus incoming edges to the same
vertex have distinct labels.

Finally fix $y\in Y_m$.  If the initial vertex
$(a_0,\ldots,a_{m-2})\in V_m$ of a lift is known, then right-resolvingness
determines at most one outgoing edge with label $y_0$, hence at most one next
state and one bit $a_{m-1}$.  Iterating forward determines all nonnegative
coordinates.  Left-resolvingness determines the negative coordinates
backwards.  Therefore there is at most one lift for each possible initial
vertex, and there are $|V_m|=2^{m-1}$ initial vertices.
\end{proof}

The arithmetic content of the threshold is already visible at $m=3$. The three
output residues $r_0,r_1,r_2$ determine the tail bit $a_4$ except for one
elementary branch, and that branch is separated by the pair $(r_0,r_2)$. The
remaining bits are then recovered from the injective map $(u,v)\mapsto u+2v$
on $\{0,1\}^2$.

\begin{lemma}[Explicit three-window decoder]\label{lem:m3-decoder}
The map
\[
\Phi_{3,3}:\{0,1\}^{5}\longrightarrow X_3^3
\]
is injective.
\end{lemma}

\begin{proof}
Write the raw block as $a_0a_1a_2a_3a_4$ and let
\[
M:=F_5=5,
\qquad
r_t:=\Val_3\bigl(\Fold_3(a_ta_{t+1}a_{t+2})\bigr)\in\{0,1,2,3,4\},
\qquad 0\le t\le 2.
\]
Because $\Val_3|_{X_3}$ is a bijection, the stabilized output triple is
equivalent to the residue triple $(r_0,r_1,r_2)$.

Using
\[
r_t\equiv a_t+2a_{t+1}+3a_{t+2}\pmod 5,
\]
set
\[
q:=r_0-r_1-r_2\pmod 5.
\]
The Fibonacci recurrence gives
\[
q\equiv a_0+a_1+2a_4\pmod 5.
\]
Since $a_0+a_1\in\{0,1,2\}$ and $2a_4\in\{0,2\}$, one has:
\[
q\in\{0,1\}\Longrightarrow a_4=0,
\qquad
q\in\{3,4\}\Longrightarrow a_4=1.
\]
Only the case $q=2$ requires an additional branch distinction. Then the two
possible triples $(a_0,a_1,a_4)$ are
\[
(1,1,0)\qquad\text{or}\qquad (0,0,1).
\]
The corresponding diagnostic residues are summarized in
Table~\ref{tab:m3-decoder}.

\begin{table}[t]
\centering
\caption{Clean branch table for the exceptional case $q=2$ in Lemma~\ref{lem:m3-decoder}.}
\label{tab:m3-decoder}
\begin{tabular}{cc}
\toprule
case for $(r_0,r_2)$ under $q=2$ & decoded value of $a_4$ \\
\midrule
$r_0=1$ & $0$ \\
$r_0=0$ & $1$ \\
$r_0=3$ and $r_2\in\{0,2\}$ & $0$ \\
$r_0=3$ and $r_2\in\{1,4\}$ & $1$ \\
\bottomrule
\end{tabular}
\end{table}

Indeed, the only possible triples are $(a_0,a_1,a_4)=(1,1,0)$ and
$(0,0,1)$. In the first branch one has
\[
r_0\equiv 3+3a_2\pmod 5,
\qquad
r_2\equiv a_2+2a_3\pmod 5,
\]
whereas in the second branch
\[
r_0\equiv 3a_2\pmod 5,
\qquad
r_2\equiv a_2+2a_3+3\pmod 5.
\]
Thus $(q,r_0,r_2)$ determines $a_4$ uniquely. Outside $q=2$ this follows from
the disjoint residue ranges above. Inside $q=2$, the values $r_0=1$ and
$r_0=0$ identify the first and second branch respectively, while the only
common value $r_0=3$ is separated by the disjoint alternatives
$r_2\in\{0,2\}$ and $r_2\in\{1,4\}$, as shown in
Table~\ref{tab:m3-decoder}.

Once $a_4$ is known, define
\[
t:=r_2-3a_4\pmod 5.
\]
Then
\[
t\equiv a_2+2a_3\pmod 5.
\]
Because the map
\[
\{0,1\}^2\to\{0,1,2,3\},\qquad (u,v)\mapsto u+2v
\]
is injective, $(a_2,a_3)$ is uniquely determined by $t$. Finally define
\[
u:=r_0-3a_2\pmod 5.
\]
Then
\[
u\equiv a_0+2a_1\pmod 5,
\]
and the same injectivity recovers $(a_0,a_1)$. Hence
$a_0a_1a_2a_3a_4$ is uniquely determined by $(r_0,r_1,r_2)$, proving
injectivity.
\end{proof}

\begin{theorem}[Sharp three-window threshold]\label{thm:block-invertibility}
Let $m\ge3$ and $n\ge m$.  Then the finite-block stabilized map
\[
\Phi_{m,n}:\{0,1\}^{n+m-1}\longrightarrow X_m^n
\]
is injective.
\end{theorem}

\begin{proof}
It is enough first to prove injectivity for $n=m$.  The case $m=3$ is
Lemma~\ref{lem:m3-decoder}.

Assume now that $m\ge4$ and that the stabilized length-$m$ output residues
\[
r_t\in\{0,\ldots,F_{m+2}-1\},\qquad 0\le t\le m-1,
\]
come from a raw block
\[
a_0a_1\ldots a_{2m-2}\in\{0,1\}^{2m-1}.
\]
For $0\le t\le m-3$, set
\[
q_t:=r_t-r_{t+1}-r_{t+2}\pmod {F_{m+2}},
\]
represented in $\{0,\ldots,F_{m+2}-1\}$.  Since
\[
r_t\equiv\sum_{j=0}^{m-1}a_{t+j}F_{j+2}\pmod {F_{m+2}},
\]
the coefficient of $a_{t+i}$ in $q_t=r_t-r_{t+1}-r_{t+2}$ is
\[
\begin{array}{c|c}
i & \text{coefficient of }a_{t+i}\\
\hline
0 & F_2=1\\
1 & F_3-F_2=1\\
2\le i\le m-1 & F_{i+2}-F_{i+1}-F_i=0\\
m & -F_{m+1}-F_m=-F_{m+2}\equiv0\\
m+1 & -F_{m+1}\equiv F_m
\end{array}
\qquad\pmod {F_{m+2}}.
\]
The interior terms cancel by the Fibonacci recurrence, and the boundary
coefficient of $a_{t+m}$ cancels because
$-F_{m+1}-F_m=-F_{m+2}$. Thus
\begin{equation}\label{eq:block-decoder-qt}
q_t\equiv a_t+a_{t+1}+F_m a_{t+m+1}\pmod {F_{m+2}}.
\end{equation}
Explicitly, the last coefficient is
$-F_{m+1}\equiv F_m\pmod {F_{m+2}}$, since
$F_{m+2}=F_{m+1}+F_m$.
If $a_{t+m+1}=0$, the right-hand side is one of $\{0,1,2\}$.  If
$a_{t+m+1}=1$, it is one of
\[
\{F_m+s:\ s=0,1,2\}.
\]
For $m\ge4$ these alternatives are disjoint, and the second set still lies
strictly below the modulus because
\[
(F_m)+2 < F_m+F_{m+1}=F_{m+2}.
\]
Thus each $q_t$ determines the boundary bit $a_{t+m+1}$.  As $t$ runs from
$0$ to $m-3$, this recovers
\[
a_{m+1},a_{m+2},\ldots,a_{2m-2}.
\]

The last residue $r_{m-1}$ now contains only two unknown bits:
\[
r_{m-1}\equiv
a_{m-1}+2a_m+\sum_{j=2}^{m-1}a_{m-1+j}F_{j+2}
\pmod {F_{m+2}}.
\]
The displayed tail sum is known.  Since
\[
(u,v)\mapsto u+2v
\]
is injective on $\{0,1\}^2$ and its values lie in $\{0,1,2,3\}$, the residue
$r_{m-1}$ recovers $(a_{m-1},a_m)$.  Then, for
$j=m-2,m-3,\ldots,0$, the residue $r_j$ has the form
\[
r_j\equiv a_j+\sum_{i=1}^{m-1}a_{j+i}F_{i+2}\pmod {F_{m+2}},
\]
where the sum is already known.  Because $a_j\in\{0,1\}$, this determines
$a_j$ uniquely.  Hence $\Phi_{m,m}$ is injective for every $m\ge4$.

Now let $n\ge m$ and suppose two raw blocks
$a_0\ldots a_{n+m-2}$ and $b_0\ldots b_{n+m-2}$ have the same
$\Phi_{m,n}$-image.  For every $0\le s\le n-m$, the length-$m$ output
subblocks beginning at $s$ agree.  By the already proved injectivity of
$\Phi_{m,m}$, the corresponding raw subblocks
\[
a_s\ldots a_{s+2m-2}
\quad\text{and}\quad
b_s\ldots b_{s+2m-2}
\]
are equal.  These intervals cover all indices $0,\ldots,n+m-2$, so the two
raw blocks are equal.
\end{proof}

\begin{corollary}[Decoder and dynamical consequences above threshold]
\label{cor:explicit-decoder}
Let $m\ge3$.  The reconstruction in Theorem
\ref{thm:block-invertibility} defines a finite map
\[
\psi_m:\mathcal L_m(Y_m)\longrightarrow\{0,1\}
\]
such that, for $y=\Phi_m(a)$,
\[
a_t=\psi_m(y_t,\ldots,y_{t+m-1}).
\]
Thus the one- and two-sided maps are conjugacies onto their images with a
memory-zero inverse of anticipation $m-1$.  Moreover, for $n\ge m$,
\[
|\mathcal L_n(Y_m)|=2^{n+m-1},
\]
and $Y_m$ is an SFT, defined for example by the allowed length-$m$ output
blocks together with the overlap compatibility of their decoded raw blocks.
Consequently
\[
h_{\rm top}(Y_m)=\log2,\qquad
\#\operatorname{Fix}(\sigma^n|_{Y_m})=2^n,
\qquad \zeta_{Y_m}(z)=\frac1{1-2z}.
\]
\end{corollary}

\begin{proof}
Apply the block reconstruction at each coordinate.  Injectivity gives the
block count, consecutive decoded blocks give a finite forbidden-word
description, and conjugacy with the full two-shift gives the entropy,
periodic-point, and zeta identities.
\end{proof}

\begin{proposition}[Sharp low-window counterexamples]\label{prop:low-window-counterexamples}
The lower-window analogues of Theorem~\ref{thm:block-invertibility} fail in
both possible senses:
\begin{enumerate}[label=(\roman*),leftmargin=*,itemsep=2pt]
  \item At $m=2$,
  \[
  \Phi_2(0^{\ZZ})=\Phi_2(1^{\ZZ})=(\ldots,00,00,00,\ldots).
  \]
  Hence the two-window sliding stabilized map is not injective.
  \item At the sharp threshold $m=3$, using only two consecutive output
  windows is not enough: the two raw blocks
  \[
  0001,\qquad 0110
  \]
  have the same $\Phi_{3,2}$-image.  Therefore the hypothesis
  $n\ge m$ in Theorem~\ref{thm:block-invertibility} cannot be replaced by
  $n\ge m-1$.
\end{enumerate}
\end{proposition}

\begin{proof}
For $m=2$ the modulus is $F_4=3$ and the two weights are
$F_2=1$ and $F_3=2$.  Thus
\[
\Val_2(00)=0,\qquad \Val_2(11)=1+2=3\equiv0\pmod 3.
\]
The unique element of $X_2$ with residue $0$ is $00$, so
\[
\Fold_2(00)=\Fold_2(11)=00.
\]
Applying this equality at every coordinate gives
\[
\Phi_2(0^{\ZZ})=\Phi_2(1^{\ZZ})=(\ldots,00,00,00,\ldots),
\]
which proves (i).

For the finite-block obstruction at $m=3$, the modulus is $F_5=5$ and the
weights are $F_2,F_3,F_4=1,2,3$.  The two consecutive windows in the raw block
$0001$ are $000$ and $001$, whose residues are
\[
0,\qquad 3.
\]
Both words are already legal, so their stabilized outputs are $000$ and
$001$.  The two consecutive windows in the raw block $0110$ are $011$ and
$110$, whose residues are
\[
2+3=5\equiv0\pmod5,\qquad 1+2=3\pmod5.
\]
The unique elements of $X_3$ with residues $0$ and $3$ are again $000$ and
$001$.  Hence
\[
\Phi_{3,2}(0001)=(000,001)=\Phi_{3,2}(0110),
\]
although $0001\ne0110$.  This proves (ii).
\end{proof}
